\paragraph{Neural TPPs for event streams.}
Neural MTPPs learn history-dependent intensities with recurrent
\citep{du2016recurrent,Mei2017TheProcess}, attention-based
\citep{zhang2020self,zuo2020thp}, jump-diffusion \citep{zhang2024njdtpp}, and
continuous-time state-space \citep{chang2025ctssm} encoders; see
\citet{Shchur2021NeuralReview} for a survey. \citet{Shi2022StateSimulation}
specialize neural Hawkes models to LOB event streams and evaluate both prediction
and simulation. Notably, the state-space line is evaluated almost exclusively
teacher-forced: \citet{chang2025ctssm} report state-of-the-art held-out likelihood
but never roll the model out (they avoid thinning altogether). Our contribution is
orthogonal to the encoder line --- we keep the state-space backbone verbatim and
change what is read \emph{out} of it, showing that the heads, not the backbone, own
closed-loop stability and calibratability.

\paragraph{Hawkes modeling of order flow.}
Parametric Hawkes processes are the classical LOB event model
\citep{Hawkes1971,Hawkes2018HawkesReview,Jain2024LimitProcess,jain2024review}, with
state-dependent extensions \citep{Morariu-Patrichi2022State-dependentModelling}.
Their linear structure gives exact stationarity and branching-ratio certificates but
limited expressiveness. \ssppp{} takes the complementary route: full neural
expressiveness with stability restored by a closed-form intensity bound rather than
by linearity.

\paragraph{Generative LOB simulators.}
Learned order-flow generators span GANs \citep{li2020stockgan,coletta2021lobgan,
coletta2022worldagent}, token-level autoregressive state-space models
\citep{nagy2023lobs5}, diffusion-guided agents \citep{huang2026diga}, and
foundation-model-scale transformers \citep{mars2025,sood2026tradefm}; agent-based
simulators \citep{byrd2020abides,frey2023jaxlob} remain the RL workhorse. These
systems are scored on distributional realism --- stylized facts
\citep{cont2001empirical,vyetrenko2020getreal} and, increasingly, the KS/Wasserstein
suites of LOB-Bench \citep{nagy2025lobbench}. Two properties separate \ssppp{} from
this line: a single likelihood-based model is held to \emph{both} axes (per-event
predictive NLL and closed-loop realism), and the simulated rate is not merely
measured but \emph{calibrated and verified}, with failures reported as first-class
results. None of the generative simulators above carries a closed-loop stability
certificate; TradeFM's inter-arrival marginal is its weakest statistic
\citep{sood2026tradefm}, and LOB-Bench's authors observe error accumulation over long
autoregressive roll-outs \citep{nagy2025lobbench} --- precisely the regime our
bounded head and rate verification target.

\paragraph{Bounded intensities.}
Bounding or normalizing intensities appears in the TPP literature mainly as a
computational device (dominating rates for thinning). Here the bound is a modeling
decision with an asymmetry requirement --- a hard ceiling for closed-loop stability,
an exactly-zero floor for quiet-regime likelihood --- and the softmin cap realizes
precisely this pair. Our cross-asset calibration study makes the case empirical:
unbounded softplus heads become progressively uncalibratable as the event rate
grows, while the bounded head admits a verified constant-rate rescaling almost
everywhere.
